\section{Training data and design of experiments}
\label{sec:doe}

The choice of training data, sampling design and fidelity strategy is
not a preprocessing step but an integral part of the surrogate model.
A design of experiments (DoE) specifies which input
configurations are evaluated---and, in adaptive settings, in what
order and at which model fidelity---to build or refine the
surrogate. For multi-energy systems (MES) in particular, that design
space spans coupled electricity, heat and sometimes cooling or gas
inputs, so the sampling plan must cover the operating region the
optimizer will explore rather than a single carrier in isolation.
For a surrogate that is only deployed in forecasting, the training
data can be drawn from historical operation and the usual forecasting
practice applies. For a surrogate that is queried repeatedly inside an
optimizer, the sampled inputs must instead cover the regions that the
optimizer will actually visit, the surrogate outputs must match the
objectives and constraints of the host problem, and the sampling
protocol must be transparent and reproducible
\cite{Perera2019191,Ruan2021221}. Reviews of learning-assisted
power-system optimization and surrogate-assisted decision making echo
the same requirement for problem-consistent simulation outputs and
explicit data-generation protocols
\cite{Khaloie2025,Lim2025,Ruan2021221}.
This section reviews the four design choices that recur most often in
the literature -- data sources, static space-filling designs, adaptive
sampling, and multi-fidelity training -- and closes with the
energy-data-specific pitfalls that affect each of them.

% BEGIN inlined table table_T3_training_doe
\begin{table*}[!t]
  \centering
  \footnotesize
  \caption{Training data and design-of-experiments practice for
  surrogate-assisted multi-energy-system and wider energy-system
  optimization. ``Pairs naturally with'' lists surrogate families most
  often combined with the strategy in the screened pool; ``Typical task''
  lists the host optimization where the strategy is most often reported.}
  \label{tab:T3-doe}
  \renewcommand{\arraystretch}{1.2}
  \begin{tabularx}{\textwidth}{@{}L{0.18\textwidth} L{0.30\textwidth} L{0.20\textwidth} L{0.26\textwidth}@{}}
    \toprule
    Strategy & Description & Pairs naturally with & Typical task (refs.) \\
    \midrule
    Latin hypercube sampling
      & One sample per input stratum for even hyper-cube coverage before metamodel fitting
      & Kriging, RBF, GP
      & Multi-carrier microgrid and DES planning, solar district heating, dispatch emulation
        \cite{Liu201914569,Seele20232480,Summ2025,Zhang2025__7,Hu2021671} \\
    Sobol / quasi-MC
      & Structured uncertainty points for polynomial-chaos propagation in stochastic optimization
      & PCE, sparse grid
      & Integrated energy dispatch and security regions; stochastic dispatch and chance-constrained OPF
        \cite{Dong2023,Jiang2026620,Gaitanis2026,Safta20172535,Lu201991827} \\
    Factorial / fractional FD
      & All or reduced combinations of discrete factor levels for response-surface fitting
      & RSM, GPR, low-order PCE
      & CCHP and integrated energy system design, hydrogen and geothermal planning
        \cite{Salahinezhad2025,Liu2025__4,Roy2023,Ahmad202451,Schulte2020} \\
    Adaptive sampling
      & Adds expensive model runs during search via infill or model-management rules
      & Kriging, RBF, GP
      & Multi-carrier microgrid planning, fuel-cell calibration, building energy design
        \cite{Liu201914569,Zhang2025__7,Aghaei Pour2022303,Granacher20221573,Lin2023} \\
    Active learning (BO)
      & Acquisition function or active-learning rule selects the next high-fidelity run
      & GP (preferred), kNN, ensembles
      & Integrated energy systems, geothermal design, energy-market assessment
        \cite{Griffith2025,Chen2025244,Chakrabarty2021,Uhrich2026211,Matta2026} \\
    Multi-fidelity training
      & Combines cheap low-fidelity sweeps with sparse high-fidelity evaluations
      & Co-kriging, multi-fidelity GNN, MF GP
      & Geothermal flow control, nuclear--renewable hybrid plants
        \cite{Chen2025244,Chung2023} \\
    Transfer learning
      & Reuses a trained surrogate in a new context with few new runs
      & NN, GBDT, GP
      & Battery state-of-charge across ambient conditions
        \cite{Chen202514792} \\
    Counterfactual / synthetic
      & Objectives and constraints from repeated high-fidelity model runs at prescribed inputs
      & NN / L2O proxies, RSM, GP
      & SCED/SCUC and OPF proxies; MES, DH and process superstructure metamodels
        \cite{Liu20223726,Wu2022231,Liu201914569,Sánchez-Zabala2024,Granacher20221573} \\
    Historical operational data
      & Past measurements; serial correlation and future information leakage if not handled
      & NN, GBDT, GP residuals
      & Multi-carrier and single-carrier load/renewable forecasting for downstream planning
        \cite{Shi2024__2,Dong2024,Koschwitz2018134,Bianchi2020244,Eseye201991463} \\
    \bottomrule
  \end{tabularx}
\end{table*}
% END inlined table table_T3_training_doe

\subsection{Data sources: synthetic, historical, hybrid}
\label{sec:doe:sources}

We distinguish three regimes. Optimizer-synthetic data is
generated by repeatedly running a high-fidelity model at prescribed
input settings and recording the objectives, constraints and other
responses the host optimizer would query. Because those outputs come
from the same model the surrogate replaces, the approximation remains
aligned with the optimization problem. In power-system studies, Liu
et al. compress the sampling space and evaluate each candidate dispatch
with full small-signal stability analysis before learning an explicit
OPF constraint \cite{Liu20223726}. Xu et al. draw dependent uncertainty
scenarios with copulas and build a polynomial-chaos response surface
from repeated AC optimal power flow runs \cite{Xu20212696}. Chen et al.
repeatedly evaluate a power-flow simulator at varied load and renewable
settings to build a training set, then fit quantile predictors for
chance-constrained OPF without explicit network parameters
\cite{Chen2023657}. For market clearing, Wu et al., Chen et al. and
Liang et al. train proxies from large batches of
security-constrained unit commitment or economic dispatch solutions
under varying load, renewable output and cost data
\cite{Wu2022231,Chen2022,Liang20246508}. For buildings, Thrampoulidis
et al. run many energy simulations for alternative retrofit and
insulation choices in Zurich and train a surrogate from the resulting
costs and system selections \cite{Thrampoulidis2021}. Aghaei Pour et al.
apply the same idea to a site with several buildings: a fast surrogate
narrows the search in repeated rounds, and the detailed building
simulation is rerun only for configurations on the current shortlist or
needed to keep the surrogate accurate \cite{Aghaei Pour2022303}. At district scale, Sánchez-Zabala
et al. combine repeated building-performance simulation with a Modelica
district-heating model to train metamodels for district optimization
\cite{Sánchez-Zabala2024}, and Chaudhry et al. evaluate a
high-fidelity fifth-generation district heating and cooling case study
over uncertain markets and demands to train a surrogate inside robust
design optimization \cite{Chaudhry2024}. Multi-carrier microgrid
and combined cooling, heating and power studies follow the same pattern
at system level: Liu et al. Latin-hypercube sample multi-energy-type
microgrid configurations and refine a kriging model with adaptive infill
\cite{Liu201914569}; Hirvonen et al. replace repeated solar-community
simulations with a neural metamodel during multi-objective sizing
\cite{Hirvonen2017323}; Gao et al. sweep battery, electrolyser and solid
oxide fuel cell sizes in MATLAB--TRNSYS co-simulation
\cite{Gao2025}. Geothermal design studies build response surfaces or
classifiers from reservoir or aquifer simulation banks---Schulte et al.
across an ensemble of geological models \cite{Schulte2020}, Petrova
et al. for well-placement screening \cite{Petrova2025}. Where a cheap
and an expensive model of the same plant both exist, Chen et al.\ combine
coarse and fine surrogates with active learning for enhanced geothermal
flow control \cite{Chen2025244}, and Chung and Wang train a co-kriging
model from a detailed Simulink representation and a fast forward
calculation of a nuclear--renewable hybrid plant linked to an industrial
process \cite{Chung2023}. Nonlinear process and utility superstructure work
selects simulation points adaptively within genetic or active-learning
loops \cite{Granacher20221573,Amoedo2023929}.
Historical operational data, by contrast, records how a system
actually ran in the past rather than how a simulation would respond to a
new design. The screened pool typically fits forecast models directly to
measured time series that downstream planning models later consume.
Eseye et al. select relevant inputs from two years of hourly building
electricity use in Otaniemi and forecast short-term demand with a neural
network \cite{Eseye201991463}. Faraji et al. train hybrid models on
measured microgrid load and weather series to support day-ahead
self-scheduling \cite{Faraji2020157284}. Wang et al. predict renewable
output and load from operating records and pass the forecasts into
hybrid-microgrid dispatch \cite{Wang20192635}. Sideratos et al. cluster
historical load trajectories and train an ensemble network for week-ahead
hourly forecasts \cite{Sideratos2020}. For coupled multi-carrier
systems, Shi and Teh forecast several load types in a regional
integrated energy system by splitting each measured history into quick
fluctuations, slower variations and a long-term trend, predicting each
part separately and adding the results back together
\cite{Shi2024__2}. For a campus with coupled electricity, heat and other
energy flows, Dong et al. update renewable forecasts each time new
measurements of generation and demand arrive, and use the latest
forecast to plan operation for the following day \cite{Dong2024}. At
district scale, Koschwitz et al. predict monthly
heating and cooling loads for a non-residential German district from
archived building consumption \cite{Koschwitz2018134}, while Bianchi
et al. extend autoregressive models of network heating load in Italy for
horizons up to 48~h \cite{Bianchi2020244}. At single-carrier sites,
Madhukumar et al. forecast day-ahead campus electricity use from several
years of hourly meter readings together with weather data
\cite{Madhukumar20228891}; Tian et al. decompose measured wind-speed
series before short-horizon prediction \cite{Tian2020}; Hanifi et al. tune
deep-learning models on recorded wind-power data \cite{Hanifi2024}; and
Li et al. forecast photovoltaic output from past generation alone for
horizons of 15--90~min \cite{Li2019}.
Hybrid approaches combine the two sources above: model-based
evaluation and recorded operation. Simulations are used to compare
configurations or operating modes that the available operating record
does not yet cover, while measurements or operating archives keep
parameters realistic and tie predictions to how the plant has actually
run. Sánchez-Zabala et al. train building
metamodels from repeated performance simulation and link them to a
Modelica district-heating network for district-scale optimisation
\cite{Sánchez-Zabala2024}. Reich et al. fit approximations from selected
recent measurements in a local heating network and embed them in dynamic
simulation for predictive operation \cite{Reich2020}. Du et al. build
district-heating surrogates from simulation runs under different thermal-
storage control strategies for a system coupled to a data centre
\cite{Du2025}. Ren et al. predict demand-response performance at
building-cluster level from operating data to size battery and thermal
storage across electricity and cooling services \cite{Ren2024}.
Jalving et al. embed market surrogates trained from a library of annual
power-market simulations into integrated energy-system design
optimisation \cite{Jalving2023}. Dong et al. fit a polynomial-chaos
model from historical integrated energy-system operation and use it in
chance-constrained day-ahead dispatch \cite{Dong2023}. Yin et al. combine
port operating records with a metamodel to size electricity and hydrogen
storage in a coupled transport--energy system \cite{Yin2024}.
The choice of regime determines the validity domain of the surrogate
(Section~\ref{sec:val:ood}) and the type of validation evidence that
can be produced (Section~\ref{sec:validation}).

\subsection{Static designs: LHS, Sobol, factorial}
\label{sec:doe:static}

When the optimum is unknown in advance, the first training set must
explore the full planning box. Static experimental designs fix
that set before any surrogate is fitted. They differ in how points are
spread through the input space and in what surrogate family they are
meant to support. Three families appear repeatedly in the screened
pool: Latin hypercube sampling (LHS), quasi-Monte Carlo and
sparse-grid collocation, and factorial or response-surface designs.

\paragraph{Latin hypercube sampling.}
LHS divides each input dimension into equal intervals and places exactly
one sample in each interval, with random permutations across dimensions.
Compared with unstructured random sampling, LHS gives more even coverage
of the planning hyper-cube; unlike factorial designs it does not require
discrete factor levels (a fixed set of values per input, such as
low/medium/high, rather than any point in a continuous range) and keeps
the number of runs manageable in moderate dimension. It is therefore widely used to initialise smooth
metamodels such as kriging or Gaussian processes before optimisation or
adaptive refinement (Section~\ref{sec:doe:adaptive}).

For a multi-carrier micro energy grid, Liu et al.\ use Latin hypercube
samples to fit a kriging model of annual cost \cite{Liu201914569}.
Seele et al.\ solve mixed-integer linear programs only at Latin hypercube
samples to calibrate linear regression surrogates of total annualized
cost for distributed energy-system planning under long-term uncertainty,
accelerating subsequent Monte Carlo design comparison
\cite{Seele20232480}. Summ et al.\ draw Latin hypercube samples from a
validated thermo-hydraulic collector model to train supervised
machine-learning metamodels before multi-objective optimisation of
insulating-glass flat-plate collectors for solar district heating
\cite{Summ2025}. Zhang et al.\ compare Latin hypercube against full
factorial sampling when calibrating a low-pressure fuel cell energy
system with a response surface and report better efficiency for LHS at
equal metamodel accuracy \cite{Zhang2025__7}. In single-carrier
dispatch emulation, Hu et al.\ apply LHS in a reduced wind-uncertainty
space to train a Gaussian-process emulator of stochastic economic
dispatch \cite{Hu2021671}.

\paragraph{Quasi-Monte Carlo, sparse grids, and polynomial-chaos collocation.}
When renewable generation, loads, or prices are uncertain, a stochastic
optimizer often needs expected costs or violation probabilities.
The straightforward Monte Carlo approach is to draw many random values
of the uncertain inputs, solve the energy model once for each draw, and
average the results. That becomes expensive when each solve is a full
dispatch or network calculation. Quasi-Monte Carlo methods replace
purely random draws with low-discrepancy point sets (Sobol and Halton
sequences are common examples) that cover the uncertainty space more
evenly. In the screened pool, however, this idea more often appears as
sparse grids or quadrature/collocation nodes used to build a
polynomial-chaos expansion (PCE): each node is one deterministic
simulator run, and the PCE coefficients summarize how outputs respond
to uncertain inputs.

Latin hypercube sampling in the previous paragraph spreads design
or planning inputs across a box to train a surrogate of cost or
performance before optimisation. Sparse-grid and PCE designs instead
pick uncertainty scenarios so the optimizer can estimate expected
costs, output spread, or violation probabilities inside a stochastic
formulation without thousands of repeated full-order solves.

For integrated energy systems, Dong et al.\ use a data-driven sparse PCE
to approximate day-ahead dispatch responses under chance constraints,
learning from historical operating data so that iterative verification
does not require repeated full-order solves \cite{Dong2023}. Jiang
et al.\ apply PCE with Galerkin projection to approximate the
security-region boundary of an electricity--gas integrated energy system,
replacing repeated nonlinear solves during boundary search
\cite{Jiang2026620}. Gaitanis et al.\ quantify capacity uncertainty in
hybrid renewable systems with combined heat and power for commercial
buildings using PCE after design optimization \cite{Gaitanis2026}. In
power-system studies, Safta et al.\ propagate wind uncertainty in
stochastic economic dispatch with PCE and sparse quadrature instead of
Monte Carlo sampling \cite{Safta20172535}; Lu et al.\ use nested
sparse-grid collocation under wind and photovoltaic uncertainty in
stochastic economic dispatch \cite{Lu201991827}; Muhlpfordt et al.\
formulate chance-constrained AC optimal power flow with PCE
\cite{Muhlpfordt20194806}; Engelmann et al.\ combine PCE with
distributed stochastic AC optimal power flow \cite{Engelmann20186188};
Koirala et al.\ translate input uncertainty to photovoltaic
hosting-capacity limits in low-voltage networks with general polynomial
chaos, avoiding repeated power-flow solves \cite{Koirala20242284}; and
Xu et al.\ use a polynomial-chaos response surface with dependent
uncertain inputs for chance-constrained AC optimal power flow
\cite{Xu20212696}. None of these papers name Sobol sequences explicitly;
they illustrate the sparse-grid, quadrature, and PCE practice that
Table~\ref{tab:T3-doe} groups under quasi-MC sampling.

\paragraph{Factorial and response-surface designs.}
Factorial and response-surface designs suit problems with only a few
settings the planner can change---for example equipment size, stack
temperature, or cooling technology---and where it matters how those
settings combine: does a larger unit still help when the operating
temperature is already high? A full factorial design runs the simulator
at every combination of discrete factor levels (a fixed set of values
per input, such as low/medium/high, rather than any point in a
continuous range)---for example all four pairings in a $2^2$ layout of
two binary factors. Fractional factorial, Box--Behnken, and related
layouts reduce the run count while still estimating main effects and
selected interactions. The fitted surface is typically a low-order
polynomial metamodel. Unlike Latin hypercube sampling, the inputs are
chosen from those predefined levels rather than spread continuously
through a box; unlike sparse-grid or PCE designs, the goal is to
explore a small set of design parameters rather than to propagate input
uncertainty through a stochastic formulation.

Salahinezhad et al.\ run a time-stepping plant model over the operating
period for each candidate CCHP design, record energy, cost and
emission indicators, and fit a response surface to those outputs before
optimising trigeneration layout with solar and hydrogen subsystems;
collector area, fuel-cell capacity, collector type and cooling
technology enter as numeric and categorical factors
\cite{Salahinezhad2025}.
Liu et al.\ use Box--Behnken designs and response-surface fits to study
how equipment capacity, price and energy coefficients affect annual cost
and carbon emissions in an integrated energy system with an energy-hub
model \cite{Liu2025__4}. Roy et al.\ employ response surface
methodology to multi-objective optimisation of a biomass combined heat
and power train with molten-carbonate fuel cell and externally fired gas
turbine, linking exergy efficiency and levelised cost to current density,
stack temperature and pressure ratio \cite{Roy2023}. Ahmad et al.\ use a
Box--Behnken design over catalyst amount, electrode voltage and
electrolysis time to model green-hydrogen production from wastewater
electrolysis before genetic and swarm optimisers search the surface
\cite{Ahmad202451}. Schulte et al.\ build Gaussian-process response
surfaces from experimental-design sample batches over 18 geological
realisations when optimising deep geothermal well placement under
uncertainty \cite{Schulte2020}. Cao et al.\ apply experimental design with
response-surface regression to select preplanned microgrid islanding
schemes under stochastic short-term simulation \cite{Cao2017}. El
Mestari et al.\ use a full factorial $2^2$ design to tune genetic-algorithm
crossover and mutation coefficients for wind-farm layout
\cite{El Mestari2025}, and Deodhar et al.\ combine $G$-optimal experimental
designs with physical tests in iterative plant-and-controller optimisation
for airborne wind energy \cite{Deodhar2018}. Priesmann et al.\ review how
design of experiments can support resource-adequacy assessment but do not
report a new factorial study themselves \cite{Priesmann2024}.

Uniform static plans---whether hypercube-filled, quadrature-based or
factorial---can still under-represent peak loads, near-feasibility bands
and active constraint boundaries. That limitation motivates the adaptive
designs in Section~\ref{sec:doe:adaptive}.

\subsection{Adaptive sampling and active learning}
\label{sec:doe:adaptive}

Static designs from the previous subsection fix all simulation points
before the surrogate is trained and used. Adaptive methods add further
points while optimisation is already running---typically where the
surrogate is inaccurate, near a constraint limit, or where the search
still needs finer resolution. Two variants are treated below: infill
rules inside surrogate-assisted search, and Bayesian optimization or
active learning that pick the next high-fidelity run.

\subsubsection{Infill rules}
\label{sec:doe:infill}

Infill rules choose the next expensive model run from the current
surrogate state. For a multi-carrier micro energy grid, Liu et al.\
start from Latin hypercube samples and then add kriging infill points
using minimum surrogate value, trust-region and mean-squared-error
criteria until the annual-cost model is accurate enough for planning
\cite{Liu201914569}. For low-pressure fuel-cell energy systems, Zhang
et al.\ use an adaptive response surface method that shrinks the
parameter search space each iteration and reuses sample points after
each reduction, cutting calibration time compared with fixed
response-surface designs \cite{Zhang2025__7}. In building
energy-system design, Aghaei Pour et al.\ search in rounds: a few
candidate layouts are shown at a time, stated preferences among their
cost-versus-emissions trade-offs steer the next search step, and a
model-management rule picks which candidates are re-evaluated with the
expensive building model to update the surrogate
\cite{Aghaei Pour2022303}. Granacher et al.\ predict whether each
candidate layout of an industrial utility system---which units to
include, how they are connected, and how they are operated---is
Pareto-optimal, and run the detailed process model next where that
prediction is most uncertain \cite{Granacher20221573}. Lin et al.\ update kriging surrogates online
with a new infilling rule while searching combined economic and
emission dispatch on large power systems \cite{Lin2023}. Deng et al.\
embed kriging in genetic and particle-swarm solvers so that most
candidate optimal power flow evaluations use the surrogate and only
promising points trigger a full AC power flow solve
\cite{Deng2020831}.

\subsubsection{Bayesian optimization and active learning}
\label{sec:doe:bo_al}

Bayesian optimization fits a Gaussian process (or related probabilistic
surrogate). An acquisition function scores where to simulate next: it
favours inputs that look strong on the surrogate and inputs where the
surrogate is still uncertain, and the highest-scoring candidate receives
the next high-fidelity run. Griffith et al.\ apply this framework to
techno-economic sizing of an integrated plant with thermal storage---the
Natrium reactor-and-storage benchmark from Idaho National
Laboratory---and report comparable or better designs with fewer
expensive model runs than gradient descent \cite{Griffith2025}. For enhanced
geothermal flow-control design, Chen et al.\ combine multi-fidelity
surrogates with active learning that queries the most informative
simulation points during optimization \cite{Chen2025244}. Chakrabarty
et al.\ calibrate 17 coupled building and HVAC parameters with scalable
Bayesian optimization, tightening and slicing the search domain as the
cost surface is learned \cite{Chakrabarty2021}. Uhrich et al.\ use
active learning to pick training runs for electricity-market metamodels
in resource-adequacy assessment, reporting large runtime reductions
with limited training data \cite{Uhrich2026211}. Matta et al.\ use
active learning with k-nearest-neighbour and Gaussian-process
surrogates to find near-optimal islanded microgrid groupings after only
a few full evaluations \cite{Matta2026}. Nikolaidis et al.\ use
Gaussian-process Bayesian optimization for data-driven unit commitment
under volatile renewable behaviour \cite{Nikolaidis2021}. Amoedo et al.\
couple active learning with a genetic algorithm and compare nine
surrogate types on a biogas separation case \cite{Amoedo2023929}. Zhang
et al.\ train surrogates for multilevel Monte Carlo resource-adequacy
assessment with a vote-by-committee active learning rule when
pre-labelled training data are unavailable \cite{Zhang2025__6}.

Reported benefits are problem-dependent, but these studies generally
reach comparable optima or Pareto sets with fewer expensive evaluations
than a one-shot static grid \cite{Liu201914569,Matta2026,Granacher20221573}.
The trade-off is tighter coupling between surrogate, query rule and host
optimizer, which complicates reproducibility and non-standard validation
(Section~\ref{sec:val:repro}).

\subsection{Multi-fidelity training and transfer learning}
\label{sec:doe:multifidelity}

Many energy-system studies already have models at different cost and
accuracy levels: a detailed time-stepping plant or network simulation
at the top, faster linearised or reduced-order versions in the middle,
and simple analytic approximations at the bottom. Multi-fidelity
training deliberately mixes data from these levels---many cheap runs
plus a few expensive ones---so the surrogate learns where the cheap
model is trustworthy and where the expensive model must be consulted.
In the screened pool, the clearest multi-carrier examples are geothermal
reservoir flow control and nuclear--renewable hybrid plants with coupled
industrial loads \cite{Chen2025244,Chung2023}. Transfer learning solves
a different problem. The simulation model stays the same, but the
operating context changes---for example another ambient temperature,
market year or geological realisation. Instead of training a surrogate
from scratch in the new context, one reuses what it already learned and
updates it with only a small number of new runs there.

\subsubsection{Combining low- and high-fidelity runs}
\label{sec:doe:mf_combine}

For enhanced geothermal systems, Chen et al.\ fit a low-fidelity
surrogate on a subset of well-control variables and a high-fidelity
surrogate on the full set; knowledge from the coarse model guides search
for the fine model, and active learning adds reservoir simulations at the
most informative points \cite{Chen2025244}. Chung and Wang study
a nuclear--renewable hybrid plant that also feeds an industrial process:
a detailed Simulink model supplies high-fidelity training data and a
fast forward calculation supplies low-fidelity data, and a co-kriging
surrogate trained on both predicts control settings across timescales
\cite{Chung2023}. Power-system studies use the same pattern at grid
level. Taghizadeh et al.\ train a graph neural network on many
low-fidelity and fewer high-fidelity power-flow solves, cutting training
cost while keeping accuracy for uncertain operating points
\cite{Taghizadeh2024}; Khayambashi et al.\ apply an enhanced
multi-fidelity graph-neural-network surrogate to chance-constrained optimal
power flow \cite{Khayambashi202453}. Xu et al.\ propagate large
sample sets through cheap polynomial-chaos surrogates and call the full
power-system model only to correct rare high-risk tail events
\cite{Xu20203127}.

\subsubsection{Transfer learning across operating conditions}
\label{sec:doe:transfer}

Multi-fidelity training mixes cheap and expensive versions of the
same model. Transfer learning keeps one model and reuses a surrogate
already trained in one context when a similar new context appears. The goal is to avoid repeating a large initial training
campaign: most of the mapping is retained, and only a few new
simulations or measurements in the new conditions are needed to adapt
it. Chen et al.\ illustrate this for battery state-of-charge estimation:
a recurrent network with a Gaussian-process layer is first trained at
one ambient temperature and then adapted quickly when the temperature
changes \cite{Chen202514792}. The same idea is attractive for
multi-objective energy-system design when a planning study moves to a
new climate year, geological realisation or market scenario, even
though the screened example above is battery estimation rather than a
full plant model.

Both multi-fidelity training and transfer learning reduce the number of
expensive evaluations needed before optimisation or uncertainty
analysis, but they also tie surrogate quality to how well the cheap
model, the source scenario and the target scenario align---a mismatch
can bias Pareto sets or constraint satisfaction if not checked in
validation (Section~\ref{sec:validation}).

\subsection{Pitfalls specific to energy data}
\label{sec:doe:pitfalls}

Energy surrogates fail quietly when training data do not match how the
system is operated or validated. Four points from the screened
literature are worth separating at the design stage.

Du et al.\ report that variability in training datasets is often
overlooked when fitting long short-term memory surrogates for district
heating with thermal storage, even though operating strategy strongly
affects robustness under peak demand \cite{Du2025}. Faraji et al.\ use
load and weather forecasts in day-ahead microgrid scheduling
\cite{Faraji2020157284}; a review of learning-assisted power-system
decision making describes exogenous forecasts as inputs to optimisation
problems \cite{Lim2025}.
Sánchez-Zabala et al.\ note that data-driven district metamodels
generalise only as far as the quality and availability of their training
simulations allow \cite{Sánchez-Zabala2024}. Extreme heat can dominate
cooling requirements \cite{Morakinyo201973}, and probabilistic risk
studies keep the full power-system model for rare high-risk tail events
that cheap surrogates mishandle \cite{Xu20203127}. These checks belong in
validation (Section~\ref{sec:validation}), not in average offline error
alone.
